\subsection{Two residual divisor classes on $X_1(2,26)$}
We first reduce the decic problem on this curve to two rational divisor
classes. The next subsection determines their effective representatives.
\begin{proposition}\label{prop:decic226-residual}
Put $X=X_1(2,26)$. Suppose an elliptic curve over a decic field
contains $C_2\times C_{26}$. Up to level automorphisms, the trace
of its level point belongs to one of two explicitly specified
rational divisor classes of degree $10$ on $X$. Both classes have a
one-dimensional Riemann--Roch space after reduction at $5$, $7$ and
$11$.
\end{proposition}
\begin{proof}
Let $D$ be the degree-ten trace divisor in the statement.
The curve has genus $25$ and good reduction at $3$. Its Jacobian has
rank zero by \cite[Theorem~4.1]{DS}.
Its $\mathbb F_3$-gonality is $12$ by \cite[Table~1]{NV}, so its
$\Q$-gonality is at least $12$. Lemma~\ref{lem:decic-cuspidal-reduction}
excludes entirely cuspidal reduction.

We enumerate all noncuspidal closed points through degree $10$,
using the same completeness checks as in
Proposition~\ref{prop:decic238-240-250}. The counts in degrees
$1,\ldots,10$ are
\[
 0,0,0,0,36,122,288,936,2328,5382.
\]
There are $301$ possible noncuspidal patterns, with cusp complements
allowed. Four maps to $X_1(26)$ leave $32$. These maps retain either
$Q$ or $P+Q$ on $E$, or retain the image of $Q$ after quotienting by
$\langle P\rangle$ or $\langle P+13Q\rangle$.

We then compute the full Picard group of $X_{\mathbb F_3}$ with
\texttt{Proof:=true}. Its degree-zero part has invariant factors
\[
 [2,14,70,1330,3990,11970].
\]
The subgroup generated by the eighteen rational cusps has invariant factors
\[
 [2,14,266,3990,11970],
\]
and index $350$.
By \cite[Proposition~4.16 and Table~3]{DEHMZ}, this is the entire
image of $J_X(\Q)$. The finite cusps have degrees $1^{18},3^6$.
We allow every cuspidal complement; in particular we include the
classes of all six cubic cusps whenever the complement degree is
at least $3$. The full Picard test leaves three patterns:
\[
 9+1_{\rm cusp},\qquad 5+5,\qquad 10.
\]
The code gives their precise closed points. Independent Sage
arithmetic checks the complete sieve, and a separate Magma run
reconstructs every Picard image using exact principal ideal
identities in both the finite and infinite maximal orders.

We first exclude the $9+1_{\rm cusp}$ pattern. Its degree-nine
noncuspidal support is fixed by
\[
 g(P,Q)=(13Q,P+17Q).
\]
This automorphism has order $3$, since its cube acts as $-1$ on
both level points. On the saved finite-field point it acts as the
third power of Frobenius. Every rational residue cusp is globally
rational, as follows from the Tate cusp calculation.
Let $C$ be the rational global cusp lifting the degree-one
complement. The divisors $gD+C$ and $D+gC$ have the same reduction.
Their difference is rational torsion and therefore vanishes by
injectivity at $3$. Both divisors have degree $11$, below the
$\Q$-gonality of $X$. They are consequently equal. Their
noncuspidal parts give $gD=D$.

A noncuspidal geometric fixed point of $g$ must have $j=0$:
an isomorphism realizing this nontrivial order-three action requires
an elliptic-curve automorphism whose image modulo $\{\pm1\}$ has
order $3$. We show that $j=0$ cannot occur here. If $E/K$ has $j=0$
and full two-torsion, write
\[
 E:y^2=x^3-\lambda^3,
 \qquad \lambda\in K^*.
\]
The other two roots force $F=\Q(\zeta_3)\subset K$, with $[K:F]=5$.
The curve is a quadratic twist of $E_0:y^2=x^3-1$, and a rational
point of order $13$ on $E$ gives an $x$-coordinate of a point of
order $13$ on $E_0$ lying in $K$. The $13$-division polynomial of
$E_0$ over $F$ has irreducible factor degrees
\[
 6,6,72.
\]
None divides $5$, a contradiction. The factorization is checked
independently in Sage. Every support point of $D$ is a conjugate of the original level point.
Thus $g$ acts freely on its support,
which would force $3\mid\deg D=10$. This excludes the first pattern.

It remains to specify the other two classes over $\Q$. Use the
normalized universal curve
\[
 y^2=x^3+cx^2+b(1-b-c)x,\qquad P=(0,0),\quad Q=(b,b).
\]
The eighteen rational cusps are labelled $C_v$ by their distinct
valuations $v=\operatorname{ord}_{C_v}(b)$. In the order
\[
 (3,-18,1,7,2,8,-22,6,-10,-2,-6,4,5,9,12,11,10,-14),
\]
the two degree-ten divisors $L_i=\sum_v a_{i,v}C_v$ have coefficient
vectors
\[
\begin{aligned}
 (a_{1,v})={}&(16,2,1,-1,-2,0,-1,-2,-1,0,0,-2,-1,0,1,0,-1,1),\\
 (a_{2,v})={}&(19,-1,-2,0,-2,0,1,-1,0,0,-2,0,0,0,0,-3,1,0).
\end{aligned}
\]
The first class reduces to the saved $5+5$ pattern; the second to
the saved degree-ten point. Exact integer lattice arithmetic verifies
these representatives. If $D$ has the corresponding reduction,
$[D-L_i]$ is rational torsion and reduces to zero at $3$, so $D\sim L_i$.

For completeness, the cusp labels are compatible across the primes
used here. Write a primitive Tate label as
$Q=\zeta_{26}^a q^{r/26}$ and
$P=(-1)^s q^{t/2}$. Put $u=\min(r,26-r)$, and let $w$ be the cusp
width. The leading terms of the normalized coordinate give
\[
 \operatorname{ord}(b)=\frac{w}{26}
 \begin{cases}-2u,&t=0,\\u,&t=1.\end{cases}
\]
They give the eighteen distinct labels above. The leading constants
are units away from $2$ and $13$, so the labels agree in every
reduction used here.

The remaining exact Riemann--Roch computations give
\[
\begin{array}{c|ccc}
 &\mathbb F_5&\mathbb F_7&\mathbb F_{11}\\\hline
 h^0(L_1)&1&1&1\\
 h^0(L_2)&1&1&1.
\end{array}
\]
These calculations give upper bounds over $\Q$. We prove effectivity
and identify the unique representatives in
Proposition~\ref{prop:decic226-complete}.
\end{proof}
